\documentclass{article}

\newcommand{\rnum}[1]{\uppercase\expandafter{\romannumeral #1\relax}}

\usepackage[utf8]{inputenc}
\usepackage[T1]{fontenc}
\usepackage[ngerman]{babel}
\usepackage[margin=2.5cm]{geometry}
\usepackage{graphicx}
\usepackage{float}
\usepackage{multirow}
\usepackage{pdflscape}
\usepackage{array}
\usepackage{enumitem}
\usepackage{longtable}

\begin{document}

\begin{titlepage}
	{
		\centering
		\huge
		\textbf{Abgabedokument Entwurf}\\
	}
	\vspace{1cm}
	\begin{table}[h]
		\centering
		\footnotesize
		\begin{tabular}{| p{1.5cm} | p{2.5cm} | p{2cm} | p{2cm} | p{2cm} |}
			\hline
			                           & Verantwortliche Teammitglieder (Organisation, Anlegen der Dokumentenversion, Moderation, Abgabe)                                   & Anwesende während der Meetings                                                                                                     & (Online-) Beiträge zum Inhalt durch                                                                                                & Korrektur-gelesen durch                                                                                                            \\
			\hline
			Dokument Entwurf Version 1 & Paul Schröder, Eric Buthmann, Yanik Buthmann, Erik Seeck, Mathis Behrendt, David Nguyen, Adrian Gode, Niklas Krüger, Kevin Zellmer & Paul Schröder, Eric Buthmann, Yanik Buthmann, Erik Seeck, Mathis Behrendt, David Nguyen, Adrian Gode, Niklas Krüger, Kevin Zellmer & Paul Schröder, Eric Buthmann, Yanik Buthmann, Erik Seeck, Mathis Behrendt, David Nguyen, Adrian Gode, Niklas Krüger, Kevin Zellmer & Paul Schröder, Eric Buthmann, Yanik Buthmann, Erik Seeck, Mathis Behrendt, David Nguyen, Adrian Gode, Niklas Krüger, Kevin Zellmer \\
			\hline
		\end{tabular}
	\end{table}
\end{titlepage}

\section{Initialer Backlog}

Auf Grundlage der vorgegebenen Anforderungen haben wir eine Liste von verfeinerten Anforderungen und User-Stories erstellt.
Zu finden ist das initiale Backlog im Verzeichnis \textit{anforderungen} im Dokument \linebreak \textit{anforderungen.pdf}.
Damit im weiteren Verlauf nachvollzogen werden kann, woher Anforderungen kommen, haben wir uns dazu entschieden jeder verfeinerten Anforderung und jeder User-Story eine ID zu geben.
Verfeinerte Anforderungen haben IDs der Form \textit{xx.x}, wobei der erste Teil die ID der ursprünglichen Anforderung und der zweite Teil eine laufende Nummer für die resultierenden Anforderungen ist.
User-Stories haben IDs der Form \textit{Sxx.x}.
Diese ID ist eine Kombination aus einem Prefix für eine User-Story und der ID der verfeinerten Anforderung.\\
Im Folgenden sollen die Verfeinerungen motiviert werden.\\
Zunächst haben wir die Benennung der Akteure in den Anforderungen vereinheitlicht.
Jeder Akteur, der das System nutzen kann, ist ein Nutzer.
Das heißt konkret, wenn in den ursprünglichen Anforderungen von Team, Teammitglied, Nutzer etc. gesprochen wird, dann sind das in den verfeinerten Anforderungen und User Stories Nutzer.
Die restlichen Akteure, welche Nutzer mit bestimmten zusätzlichen Eigenschaften sind (z.B. Product Owner), wurden nicht umbenannt.\\
Desweiteren ist ein großer Teil der Anforderungen nicht atomar.
Daher haben wir die meisten Anforderungen in mehrere verfeinerte Anforderungen aufgeteilt.
Zum Beispiel haben wir die erste Anforderung in drei verfeinerte Anforderungen aufgeteilt, welche die Eigenschaft der Atomarität besitzen.
Dieses Beispiel soll exemplarisch die Aufteilung der Anforderungen zeigen.
Im weiteren Verlauf wird aufgrund der Existenz eines Beispiels und der Übertragbarkeit des Beispiels auf die anderen Fälle nicht weiter auf die Eigenschaft der Atomarität eingegangen.\\
Die zweite Anforderung fordert unter anderem, dass ein Product Owner einer User-Story eine Priorität zuordnen kann.
Diese Anforderung ist nicht eindeutig.
Wir haben uns daher entschieden, den Begriff \textit{Priorität} zu konkretisieren und die Prioritätsstufen \textit{hoch}, \textit{normal} und \textit{niedrig} zu definieren.\\
Außerdem haben wir uns dafür entschieden, den anderen Teil der Anforderung umzuformulieren.
Neue User Stories werden einem Backlog hinzugefügt, da es mehrere Backlogs geben kann.\\
Die dritte Anforderung beinhaltet die Priorisierung von Tasks.
Diese Anforderung haben wir analog zur zweiten Anforderung behandelt und Prioritätsstufen definiert.\\
Anforderung 4 ist nicht eindeutig.
Der Begriff Liste wird mehrmals für unterschiedliche Dinge verwendet.
Daher haben wir uns entschieden, verschiedene Begriffe für Listen in unterschiedlichen Kontexten einzuführen.
Ein \textit{Backlog} ist eine Liste an User Stories und Tasks.
In einem \textit{Backlog} gibt es mehrere Listen an Tasks.
Diese Tasklisten beinhalten alle Tasks, die sich in einem gewissen Status befinden.
Daher haben wir Listen in diesem Kontext \textit{Tasklisten mit dazugehörigem Status} genannt.
Die Aktion \textit{Task verschieben} kann nun in unterschiedlichen Kontexten ausgeführt werden.
Daher haben wir für jeden Kontext eine separate Anforderung und User Story erstellt.
Außerdem haben wir den Zustand \textit{Wartet} hinzugefügt, da dieser unserer Meinung nach essentiell für Tasks ist, an denen noch nicht gearbeitet wird.\\
Anforderung 5 ist ebenfalls nicht eindeutig, da nicht konkret darauf eingegangen wird, welche Menge an Informationen ein Profil besitzen soll.
Wir haben daher diese Menge exakt definiert.
Zusätzlich haben wir definiert, welche Informationen in einem Profil geändert werden dürfen.
Zum Beispiel ist die Änderung der Rolle durch den Nutzer selbst ist unserer Meinung nach nicht sinnvoll.
Daher haben wir diese Information aus der Menge der editierbaren Eigenschaften ausgeschlossen.\\
Die Anforderung mit der ID 7 ist nicht einzigartig, denn diese wird bereits durch Anforderung 6 abgedeckt.
Daher haben wir daraus keine verfeinerte Anforderung entwickelt.\\
Die 8. Anforderung fordert eine Visualisierung der Abweichungen zwischen geschätztem Aufwand und tatsächlichem Aufwand.
Der Begriff \textit{Visualisierung} ist nicht eindeutig.
Daher haben wir uns entschieden die Anforderung zu konkretisieren und als Visualisierungsart ein Balkendiagramm zu wählen.\\
In Anforderung 10 werden Benachrichtigungen mittels E-Mail oder in-system gefordert.
E-Mail Benachrichtigungen sind im Kontext dieser Projektarbeit unpraktisch.
Daher haben wir uns entschieden nur in-system Benachrichtigungen zu fordern.\\
Die Anforderung mit der ID 11 haben wir entfernt, da diese nicht einzigartig ist und bereits von anderen Anforderungen abgedeckt wird, und ebenfalls nicht messbar ist.
Zum Beispiel fordert Anforderung 7 die Zuordnung von Tasks zu mehreren Nutzern, um Pair Programming zu unterstützen.
Außerdem unterstützt das System inhärent die Agile Praktik \textit{Sprint Planung}.
Wir haben entschieden, dass das der Formulierung ``so viele Agile Practices wie möglich`` genügt.\\
Anforderung 16 ist eine Konkretisierung von Anforderung 12.
Daher haben wir uns entschieden, diese Anfoderungen zu vereinigen.\\
Die Anforderungen 14 und 15 sind beide nicht eindeutig.
Im Fall der Anforderungen 14 haben wir uns dazu entschieden, 10 gleichzeitige Nutzer zu unterstützen.
Im Rahmen dieser Projektarbeit erachten wir diese Anzahl als praktisch, obwohl wir wissen, dass in der Realität deutlich höhere Nutzerzahlen zu erwarten wären.
Die 15. Anforderung haben wir dahingehend konkretisiert, dass das System eine Verfügbarkeit von 99\% erbringen soll.
Im Rahmen dieser Projektarbeit können wir diese Metrik jedoch nicht prüfen, da diese unter anderem von Faktoren wie dem Update-Zyklus abhängt.
Daher soll die genannte Zahl als fiktiv betrachtet werden.\\
Die Gesamtheit der Use Cases kann der Abbildung ... entnommen werden.\\ %TODO
Die entwickelten User Stories werden in das Gitlab übertragen.
Während der Sprint Plaung werden diese dann in ausführbare Aufgaben übersetzt.

\end{document}
